% !TeX encoding = UTF-8
\documentclass[variant=guia,cover=institutional,mode=screen]{ua-engineering-v6-article}
% !TeX encoding = UTF-8
\UASetUniversity{Universidad Austral}
\UASetFaculty{Facultad de Ingeniería}
\UASetDepartment{Departamento de Computación}
\UASetCampus{Pilar}
\UASetCareer{Ingeniería Informática}
\UASetCommission{Comisión A}
\UASetProfessor{Equipo docente}
\UASetSemester{Segundo cuatrimestre 2026}
\UASetCourse{Sistemas Distribuidos}
\UASetDate{2026}


\UASetClassNumber{10}
\UASetClassTitle{Observabilidad distribuida}
\UASetDocKind{Guía de ejercicios}

\title{Clase 10 --- Guía de ejercicios}
\author{\UAProfessor}
\date{\UADate}

\begin{document}

\section{Indicaciones generales}
Resuelva cada ejercicio declarando supuestos, modelo de fallas, garantía y trade-off. Cuando corresponda, construya una ejecución verificable y vincule la respuesta con evidencia observable. Justificar siempre. En esta clase no alcanza con decir “usar logs” o “mirar trazas”: se espera explicar qué problema resuelve cada herramienta, qué límites tiene y qué patrones de diseño exige.

\section{Nivel 1: fundamentos y reconocimiento}
\begin{uaexercise}
Defina observabilidad distribuida y explique en qué se diferencia de monitoreo básico.
\end{uaexercise}

\begin{uaexercise}
Explique al menos tres razones por las cuales un sistema distribuido en producción necesita monitoreo activo, y por qué cada una no se satisface solo con testing pre-producción.
\end{uaexercise}

\begin{uaexercise}
Un equipo recibe una alerta que dice ``error rate > 5\%''. Explique por qué esta alerta, por sí sola, es insuficiente para actuar. ¿Qué información adicional necesitaría el equipo para responder ``qué está roto'' y ``por qué está roto''?
\end{uaexercise}

\begin{uaexercise}
Explique el antipatrón de las ``Tres Pestañas del Navegador'' (\emph{Three Browser Tabs}) y compárelo con el modelo de ``La Trenza Única de Datos'' (\emph{Single Braid of Data}).
\end{uaexercise}

\begin{uaexercise}
Diferencie \textbf{Hard Context} de \textbf{Soft Context}. Explique por qué el Hard Context es indispensable para que las computadoras puedan calcular correlaciones automáticas entre telemetrías.
\end{uaexercise}


\section{Nivel 2: ejecuciones y fallas}
\begin{uaexercise}
Compare logs, métricas y trazas en términos de qué pregunta responde cada uno y cómo se entrelazan mediante OpenTelemetry.
\end{uaexercise}

\begin{uaexercise}
Explique por qué los logs estructurados son preferibles a los logs de texto libre en sistemas distribuidos.
\end{uaexercise}

\begin{uaexercise}
Defina la noción de \textbf{Evento Estructurado Ancho} (\emph{Arbitrarily Wide Structured Event}) y explique por qué agrupar la telemetría de una unidad de trabajo es superior a emitir múltiples líneas de log sueltas.
\end{uaexercise}

\begin{uaexercise}
Diseñe un esquema JSON de campos para un evento estructurado en un servicio HTTP de comercio electrónico.
\end{uaexercise}

\begin{uaexercise}
Construya un escenario donde la ausencia de \texttt{trace\_id} o \texttt{span\_id} en un log impida determinar el impacto de un error en el cliente final.
\end{uaexercise}


\section{Nivel 3: diseño y comparación}
\begin{uaexercise}
Explique qué son las Golden Signals y por qué son útiles para evaluar la salud de un servicio.
\end{uaexercise}

\begin{uaexercise}
Explique qué es la \textbf{Explosión de Cardinalidad} en bases de datos de series temporales (TSDB como Prometheus) y por qué los identificadores de alta cardinalidad (\texttt{user\_id}, \texttt{order\_id}) pertenecen a trazas/eventos y no a métricas agregadas.
\end{uaexercise}

\begin{uaexercise}
Explique por qué mirar promedios de latencia puede ocultar problemas serios que sí se reflejan en el percentil p99.
\end{uaexercise}

\begin{uaexercise}
\textbf{Conexión con Clase 03:} Explique por qué el desvío de reloj físico (\emph{clock drift}) entre servidores impide ordenar eventos distantes por timestamp real y cómo el campo \texttt{parent\_span\_id} implementa la relación \emph{happened-before} ($\rightarrow$) de Lamport mediante un Grafo Acíclico Dirigido (DAG).
\end{uaexercise}

\begin{uaexercise}
Diferencie la propagación de contexto \textbf{In-Band} (en cabeceras HTTP/gRPC como W3C \texttt{traceparent}) de la recolección \textbf{Out-of-Band} (exportación asíncrona OTLP al Collector).
\end{uaexercise}


\section{Nivel 4: integración y defensa}
\begin{uaexercise}
Explique la \textbf{Estrategia de Instrumentación en 2 Capas}: Capa 1 de Infraestructura (Auto-instrumentación / esqueleto de propagación) vs. Capa 2 de Dominio (Enriquecimiento semántico manual).
\end{uaexercise}

\begin{uaexercise}
Construya un escenario de cascada de fallas donde el síntoma visible sea un error HTTP 504 en la API de borde, pero la causa raíz sea una base de datos downstream saturada.
\end{uaexercise}

\begin{uaexercise}
Describa los 4 pasos del \textbf{Core Analysis Loop} (Análisis por Primeros Principios) para diagnosticar un incidente en producción y explique por qué ``buscar la línea exacta que falló'' no es un enfoque sostenible.
\end{uaexercise}

\begin{uaexercise}
Explique cómo cambia el problema de propagación de contexto en un sistema basado en eventos (Event-Driven / Kafka) respecto de uno sincrónico (HTTP/gRPC).
\end{uaexercise}

\begin{uaexercise}
Aplique el modelo del curso $D=(P,F,G,S,T)$ al problema de la observabilidad distribuida.
\end{uaexercise}

\end{document}
